% !TeX root = ./thesis.tex
% !TeX spellcheck = hu_HU
% !TeX encoding = UTF-8
% !TeX program = pdflatex
% !TeX TXS-program:compile = txs:///python//usr/bin/python3
\documentclass[12pt,a4paper,oneside]{report}             % Egyoldalas (javasolt)
%\documentclass[11pt,a4paper,twoside,openright]{report}  % Duplex
% thanks to http://tex.stackexchange.com/a/47579/71109
\usepackage{pdfpages} 
\usepackage{ifxetex}
\usepackage{ifluatex}
\newif\ifxetexorluatex % a new conditional starts as false
\ifnum 0\ifxetex 1\fi\ifluatex 1\fi>0
   \xetexorluatextrue
\fi

\ifxetexorluatex
  \usepackage{fontspec}
\else
  \usepackage[T1]{fontenc}
  \usepackage[utf8]{inputenc}
  \usepackage[lighttt]{lmodern}
\fi

\usepackage[english,magyar]{babel} % Alapértelmezés szerint utoljára definiált nyelv lesz aktív, de később külön beállítjuk az aktív nyelvet.

%\usepackage{cmap}
\usepackage{amsfonts,amsmath,amssymb} % Mathematical symbols.
%\usepackage[ruled,boxed,resetcount,linesnumbered]{algorithm2e} % For pseudocodes. % beware: this is not compatible with LuaLaTeX, see http://tex.stackexchange.com/questions/34814/lualatex-and-algorithm2e
\usepackage{booktabs} % For publication quality tables for LaTeX
\usepackage{graphicx}

%\usepackage{fancyhdr}
%\usepackage{lastpage}

\usepackage{anysize}
%\usepackage{sectsty}
\usepackage{setspace} % For setting line spacing

\usepackage[unicode]{hyperref} % For hyperlinks in the generated document.
\usepackage{xcolor}
\usepackage{listings} % For source code snippets.

\usepackage[amsmath,thmmarks]{ntheorem} % Theorem-like environments.

\usepackage[hang]{caption}

\usepackage[list=true]{subcaption}	%tof url
\usepackage{tocloft}

\usepackage{datetime2}


\singlespacing

\newcommand{\selecthungarian}{
	\selectlanguage{magyar}
	\setlength{\parindent}{2em}
	\setlength{\parskip}{0em}
	\frenchspacing
}

\newcommand{\selectenglish}{
	\selectlanguage{english}
	\setlength{\parindent}{0em}
	\setlength{\parskip}{0.5em}
	\nonfrenchspacing
	\renewcommand{\figureautorefname}{Figure}
	\renewcommand{\tableautorefname}{Table}
	\renewcommand{\partautorefname}{Part}
	\renewcommand{\chapterautorefname}{Chapter}
	\renewcommand{\sectionautorefname}{Section}
	\renewcommand{\subsectionautorefname}{Section}
	\renewcommand{\subsubsectionautorefname}{Section}
}

\usepackage[numbers]{natbib}
\usepackage{xspace}

\usepackage{url}
\usepackage{multicol} %stuff side by side
\usepackage{ifthen}

% Ide teheted a packageket amiket használni szeretnél

%--------------------------------------------------------------------------------------
\usepackage{float}
\usepackage{hyperref}
\usepackage{url}
\usepackage{caption}
\usepackage{graphicx} % Képek kezeléséhez
\usepackage{listings} % Kódrészletekhez
\usepackage{booktabs} % Táblázatokhoz
\usepackage{comment}
\usepackage[magyar]{babel}
\usepackage{amsmath}
\usepackage{ragged2e}
\usepackage{array}
\usepackage{times}
\usepackage{fancyhdr}
\usepackage[a4paper, left=3cm, right=2.54cm, top=2.54cm, bottom=2.54cm]{geometry}
% Fejléc és lábléc
\pagestyle{fancy}
\fancyhf{}
\fancyhead[L]{\textnormal{\thetitle}}
\fancyhead[C]{\szerzoNeptun}
\fancyfoot[C]{\thepage}
\renewcommand{\headrulewidth}{0.5pt}
\renewcommand{\footrulewidth}{0pt}

\usepackage{titlesec}
% 1. szint: Decimális számozás, 3.1.3-ban meghatározott stílus
\titleformat{\section}[block]
  {\bfseries\fontsize{16pt}{19.2pt}\selectfont} % Félkövér, 16pt betűméret
  {\thesection.} % Számozás
  {1em} % Címszöveg előtti távolság
  {} % Címszöveg formázása

\titlespacing{\section}
  {0pt} % Bal oldali behúzás
  {24pt} % Cím előtti térköz
  {18pt} % Cím utáni térköz

% 2. szint: Decimális számozás, stílus szerint
\titleformat{\subsection}[block]
  {\bfseries\fontsize{14pt}{16.8pt}\selectfont} % Félkövér, 14pt betűméret
  {\thesubsection.}
  {1em}
  {}

\titlespacing{\subsection}
  {0pt}
  {18pt}
  {12pt}

% 3. szint: Decimális számozás, stílus szerint
\titleformat{\subsubsection}[block]
  {\normalfont\fontsize{14pt}{16.8pt}\selectfont} % Normál, 14pt betűméret
  {\thesubsubsection.}
  {1em}
  {}

\titlespacing{\subsubsection}
  {0pt}
  {12pt}
  {6pt}

% 4. és további szintek: egyéni meghatározás
\titleformat{\subsubsubsection}[block]
  {\normalfont\fontsize{14pt}{16.8pt}\selectfont} % Normál, 14pt betűméret
  {\thesubsubsubsection.}
  {1em}
  {}

\titlespacing{\subsubsubsection}
  {0pt}
  {12pt}
  {6pt}

% Alapértelmezett képaláírás beállítás
\captionsetup{
  font=bf, % Félkövér
  labelfont=bf,
  format=plain,
  justification=centering, % Középre igazítás
  textfont={bf,rm}, % Times New Roman betűtípus
  labelsep=period, % Pont a típus és a cím között
  %skip=6pt, % Cím előtti térköz
  %belowskip=18pt % Cím utáni térköz
}

% Képeknél és ábráknál a ”képaláírás” az objektum alá, táblázatoknál és kódrészleteknél az objektum fölé kerüljön
%\captionsetup[figure]{position=bottom}
%\captionsetup[table]{position=top}
%\captionsetup[lstlisting]{position=top}


%--------------------------------------------------------------------------------------
\newcommand{\szerzoVezeteknev}{Székely}
\newcommand{\szerzoKeresztnev}{Dániel}
\newcommand{\szerzoNeptun}{JAXC3C}

\newcommand{\szakirany}{} % Informatikusoknál nincs szakirány. Villamosmérnököknél: Automatizálás (\aut) vagy Infokommunikáció (\infokom).

\newcommand{\konzulensAMegszolitas}{}
\newcommand{\konzulensAVezeteknev}{Paál}
\newcommand{\konzulensAKeresztnev}{Dávid}
\newcommand{\konzulensBMegszolitas}{}
\newcommand{\konzulensBVezeteknev}{Tamás}
\newcommand{\konzulensBKeresztnev}{Dávid}
\newcommand{\konzulensCMegszolitas}{}
\newcommand{\konzulensCVezeteknev}{}
\newcommand{\konzulensCKeresztnev}{}

\newcommand{\cim}{Mitől leszek jó vezető? Egyéni kompetencia-elemzés és fejlesztési terv} % Cím
\newcommand{\tanszek}{\szeit} % informatika (\szeit), automatizálási (\szeaut) vagy távközlési (\szetat)
\newcommand{\szak}{\infoMSc} % Mérnökinformatikus BSc (\infoMsc), MSc (\infoMsc), Gazdaságinformatikus BSc (\gazdInfoBsc), MSc (\gazdInfoMsc), vagy Villamosmérnöki BSc (\villBSc), MSc (\villMSc)

%--------------------------------------------------------------------------------------
% Elnevezések
%--------------------------------------------------------------------------------------

\newif\ifen
\newif\ifhu

\newcommand{\en}[1]{\ifen#1\fi}
\newcommand{\hu}[1]{\ifhu#1\fi}


\newcommand{\sze}{%
    \en{Széchenyi István University}%
    \hu{Széchenyi István Egyetem}%
}
\newcommand{\givk}{%
    \en{Faculty of Mechanical Engineering, Informatics and Electrical Engineering}%
    \hu{Gépészmérnöki, Informatikai és Villamosmérnöki Kar}%
}
\newcommand{\szeit}{%
    \en{Department of Informatics}%
    \hu{Informatika Tanszék}%
}
\newcommand{\szeaut}{%
    \en{Department of Automation}%
    \hu{Automatizálási Tanszék}%
}
\newcommand{\szetat}{%
    \en{Department of Telecommunications}%
    \hu{Távközlési Tanszék}%
}
\newcommand{\aut}{%
    \en{Specialization in Automation}%
    \hu{Automatizálási Szakirány}%
}
\newcommand{\infokom}{%
    \en{Specialization in Infocommunication}%
    \hu{Infokommunikáció Szakirány}%
}
\newcommand{\keszitette}{%
    \en{Author}%
    \hu{Készítette}%
}
\newcommand{\konzulens}{%
    \en{Advisor}%
    \hu{Konzulens}%
}
\newcommand{\szakdolgozat}{%
    \en{Bachelor's Thesis}%
    \hu{Egyéni dolgozat}%
}
\newcommand{\diplomaterv}{%
    \en{Master's Thesis}%
    \hu{Diplomamunka}%
}
\newcommand{\dolgozat}{%
    \en{Project}%
    \hu{Dolgozat}%
}
\newcommand{\infoBSc}{%
    \en{Computer Science Engineering, BSc}%
    \hu{Mérnökinformatikus BSc}%
}
\newcommand{\infoMSc}{%
    \en{Computer Science Engineering, MSc}%
    \hu{Mérnökinformatikus MSc}%
}
\newcommand{\szeGyor}{%
    \en{SZÉCHENYI ISTVÁN UNIVERSITY}%
    \hu{SZÉCHENYI ISTVÁN EGYETEM}%
}
\newcommand{\gazdInfoBSc}{%
    \en{Business Informatics, BSc}%
    \hu{Gazdaságinformatikus BSc}%
}
\newcommand{\gazdInfoMSc}{%
    \en{Business Informatics, MSc}%
    \hu{Gazdaságinformatikus MSc}%
}
\newcommand{\villBSc}{%
    \en{Electrical Engineering BSc}%
    \hu{Villamosmérnöki BSc}%
}
\newcommand{\villMSc}{%
    \en{Electrical Engineering MSc}%
    \hu{Villamosmérnöki MSc}%
}
\newcommand{\pelda}{%
    \en{Example}%
    \hu{Példa}%
}
\newcommand{\definicio}{%
    \en{Definition}%
    \hu{Definíció}%
}
\newcommand{\tetel}{%
    \en{Theorem}%
    \hu{Tétel}%
}
\newcommand{\bevezetes}{%
    \en{Introduction}%
    \hu{Bevezetés}%
}
\newcommand{\fuggelek}{%
    \en{Appendix}%
    \hu{Függelék}%
}
\newcommand{\Literature}{%
    \en{Management and Leadership Theoretical Background}%
    \hu{A vezetés és kompetenciák elméleti háttere}%
}
\newcommand{\ManagementExperience}{%
    \en{Own Leadership Experience and Introduction}%
    \hu{Saját vezetői tapasztalat és bemutatkozás}%
}
\newcommand{\Introspection}{%
    \en{Self Assessment Test with the Competence Analysis}%
    \hu{Egyéni kompetenciaelemzés az önértékelő teszt alapján}%
}
\newcommand{\DevelopementPlan}{%
    \en{Development Plan}%
    \hu{Fejlesztési terv}%
}



\newcommand{\szerzo}{%
    \en{\szerzoKeresztnev{} \szerzoVezeteknev}%
    \hu{\szerzoVezeteknev{} \szerzoKeresztnev}%
}
\newcommand{\konzulensA}{%
    \en{\konzulensAMegszolitas\konzulensAKeresztnev{} \konzulensAVezeteknev}%
    \hu{\konzulensAMegszolitas\konzulensAVezeteknev{} \konzulensAKeresztnev}%
}
\newcommand{\konzulensB}{%
    \en{\konzulensBMegszolitas\konzulensBKeresztnev{} \konzulensBVezeteknev}%
    \hu{\konzulensBMegszolitas\konzulensBVezeteknev{} \konzulensBKeresztnev}%
}
\newcommand{\konzulensC}{%
    \en{\konzulensCMegszolitas\konzulensCKeresztnev{} \konzulensCVezeteknev}%
    \hu{\konzulensCMegszolitas\konzulensCVezeteknev{} \konzulensCKeresztnev}%
}

% Beállítások magyar nyelvű dolgozathoz

\hutrue

\newcommand{\selectthesislanguage}{
    \selecthungarian
    \enfalse
    \hutrue
}

\bibliographystyle{huplain}

% Opcionálisan átnevezhető címek
%\addto\captionsmagyar{%
%\renewcommand{\listfigurename}{Saját ábrajegyzék cím}
%\renewcommand{\listtablename}{Saját táblázatjegyzék cím}
%\renewcommand{\bibname}{Saját irodalomjegyzék név}
%}

\def\lstlistingname{lista}

\newcommand{\appendixnumber}{6}  % a fofejezet-szamlalo az angol ABC 6. betuje (F) lesz

% Settings for English documents
%
\entrue

\newcommand{\selectthesislanguage}{
    \selectenglish
    \hufalse
    \entrue
}

\bibliographystyle{plainnat}

% Optional custom titles
%\addto\captionsenglish{%
%\renewcommand*{\listfigurename}{Your list of figures title}
%\renewcommand*{\listtablename}{Your list of tables title}
%\renewcommand*{\bibname}{Your bibliography title}
%}

\newcommand{\ie}{i.e.\@\xspace}
\newcommand{\Ie}{I.e.\@\xspace}
\newcommand{\eg}{e.g.\@\xspace}
\newcommand{\Eg}{E.g.\@\xspace}
\newcommand{\etal}{et al.\@\xspace}
\newcommand{\etc}{etc.\@\xspace}
\newcommand{\vs}{vs.\@\xspace}
\newcommand{\viz}{viz.\@\xspace} % videlicet
\newcommand{\cf}{cf.\@\xspace} % confer
\newcommand{\Cf}{Cf.\@\xspace}
\newcommand{\wrt}{w.r.t.\@\xspace} % with respect to

\newcommand{\appendixnumber}{1}  % a fofejezet-szamlalo az angol ABC 1. betuje (A) lesz



\newcommand{\szerzoMeta}{\szerzoVezeteknev{} \szerzoKeresztnev} % egy szerző esetén TODO@FMA két szerző
\newcommand{\doktipus}{\szakdolgozat} % Dokumentum típusa (\szakdolgozat, \diplomaterv vagy \dolgozat)

%--------------------------------------------------------------------------------------
% Page layout setup
%--------------------------------------------------------------------------------------
% we need to redefine the pagestyle plain
% another possibility is to use the body of this command without \fancypagestyle
% and use \pagestyle{fancy} but in that case the special pages
% (like the ToC, the References, and the Chapter pages)remain in plane style

\pagestyle{plain}
\marginsize{35mm}{25mm}{15mm}{15mm}

\setcounter{tocdepth}{3}
%\sectionfont{\large\upshape\bfseries}
\setcounter{secnumdepth}{3}

\sloppy % Margón túllógó sorok tiltása.
\widowpenalty=10000 \clubpenalty=10000 %A fattyú- és árvasorok elkerülése
\def\hyph{-\penalty0\hskip0pt\relax} % Kötőjeles szavak elválasztásának engedélyezése


%--------------------------------------------------------------------------------------
% Setup hyperref package
%--------------------------------------------------------------------------------------
\hypersetup{
    % bookmarks=true,            % show bookmarks bar?
    unicode=true,              % non-Latin characters in Acrobat's bookmarks
    pdftitle={\cim},        % title
    pdfauthor={\szerzoMeta},    % author
    pdfsubject={\doktipus}, % subject of the document
    pdfcreator={\szerzoMeta},   % creator of the document
    pdfproducer={},    % producer of the document
    pdfkeywords={},    % list of keywords (separate then by comma)
    pdfnewwindow=true,         % links in new window
    colorlinks=true,           % false: boxed links; true: colored links
    linkcolor=black,           % color of internal links
    citecolor=black,           % color of links to bibliography
    filecolor=black,           % color of file links
    urlcolor=black             % color of external links
}


%--------------------------------------------------------------------------------------
% Set up listings
%--------------------------------------------------------------------------------------
\definecolor{lightgray}{rgb}{0.95,0.95,0.95}
\lstset{
	basicstyle=\scriptsize\ttfamily, % print whole listing small
	keywordstyle=\color{black}\bfseries, % bold black keywords
	identifierstyle=, % nothing happens
	% default behavior: comments in italic, to change use
	% commentstyle=\color{green}, % for e.g. green comments
	stringstyle=\scriptsize,
	showstringspaces=false, % no special string spaces
	aboveskip=3pt,
	belowskip=3pt,
	backgroundcolor=\color{lightgray},
	columns=flexible,
	keepspaces=true,
	escapeinside={(*@}{@*)},
	captionpos=b,
	breaklines=true,
	frame=single,
	float=!ht,
	tabsize=2,
	literate=*
		{á}{{\'a}}1	{é}{{\'e}}1	{í}{{\'i}}1	{ó}{{\'o}}1	{ö}{{\"o}}1	{ő}{{\H{o}}}1	{ú}{{\'u}}1	{ü}{{\"u}}1	{ű}{{\H{u}}}1
		{Á}{{\'A}}1	{É}{{\'E}}1	{Í}{{\'I}}1	{Ó}{{\'O}}1	{Ö}{{\"O}}1	{Ő}{{\H{O}}}1	{Ú}{{\'U}}1	{Ü}{{\"U}}1	{Ű}{{\H{U}}}1
}


%--------------------------------------------------------------------------------------
% Set up theorem-like environments
%--------------------------------------------------------------------------------------
% Using ntheorem package -- see http://www.math.washington.edu/tex-archive/macros/latex/contrib/ntheorem/ntheorem.pdf

\theoremstyle{plain}
\theoremseparator{.}
\newtheorem{example}{\pelda}

\theoremseparator{.}
%\theoremprework{\bigskip\hrule\medskip}
%\theorempostwork{\hrule\bigskip}
\theorembodyfont{\upshape}
\theoremsymbol{{\large \ensuremath{\centerdot}}}
\newtheorem{definition}{\definicio}

\theoremseparator{.}
%\theoremprework{\bigskip\hrule\medskip}
%\theorempostwork{\hrule\bigskip}
\newtheorem{theorem}{\tetel}


%--------------------------------------------------------------------------------------
% Some new commands and declarations
%--------------------------------------------------------------------------------------
\newcommand{\code}[1]{{\upshape\ttfamily\scriptsize\indent #1}}
\newcommand{\doi}[1]{DOI: \href{http://dx.doi.org/\detokenize{#1}}{\raggedright{\texttt{\detokenize{#1}}}}} % A hivatkozások közt így könnyebb DOI-t megadni.

\DeclareMathOperator*{\argmax}{arg\,max}
%\DeclareMathOperator*[1]{\floor}{arg\,max}
\DeclareMathOperator{\sign}{sgn}
\DeclareMathOperator{\rot}{rot}


%--------------------------------------------------------------------------------------
% Setup captions
%--------------------------------------------------------------------------------------
\captionsetup[figure]{
	width=.75\textwidth,
	aboveskip=10pt}

\renewcommand{\captionlabelfont}{\bf}
%\renewcommand{\captionfont}{\footnotesize\it}

%--------------------------------------------------------------------------------------
% Hyphenation exceptions
%--------------------------------------------------------------------------------------
\hyphenation{Shakes-peare Mar-seilles ár-víz-tű-rő tü-kör-fú-ró-gép}

%--------------------------------------------------------------------------------------
% Sources for figures
%--------------------------------------------------------------------------------------

% takes 3 arguments: URL, month, day
\makeatletter
\newcommand{\figsourcefont}{\footnotesize}
\newcommand{\figsource}[3]{%
  \addtocontents{lof}{%
    {\leftskip\cftfigindent
     \advance\leftskip\cftfignumwidth
     \rightskip\@tocrmarg 
     \scriptsize \url{#1} \newline Utolsó látogatás időpontja: \DTMdate{\the\year-#2-#3}
     \par}%
  }%
 }
\makeatother


\author{\szerzo}
\title{\title}


%--------------------------------------------------------------------------------------
% Useful math macros
%--------------------------------------------------------------------------------------

% Pár előre definiált makró segítségképp
\newcommand{\EqHMargin}{\\[0.1cm]}                                                                             % egyenletek közötti helykihagyás
\newcommand{\vek}[1]{\MakeLovercase{\mathbf{#1}}}                                             % vektor jelölése
\newcommand{\mat}[1]{\MakeUppercase{\mathbf{#1}}}                                          % mátrix jelölése
\newcommand{\rotacio}[1]{\nabla \times \MakeUppercase{\mathbf{#1}}}         % rotáció
\newcommand{\divergencia}[1]{\nabla \cdot \MakeUppercase{\mathbf{#1}}}  % divergencia

% Ide teheted a saját makróidat
 % beállítások, nem kell vele foglalkoznod remélhetőleg, de ha valami latex hekkelésre vagy új parancsra van szükséged annak itt a helye

% A default szöveg legyen times new roman 12-es


%--------------------------------------------------------------------------------------
% Itt kezdődik a dolgozat
%--------------------------------------------------------------------------------------
% Másfeles sorköz
\newcommand{\masfelessorkoz}{\renewcommand{\baselinestretch}{1.24}\small\normalsize}
%--------------------------------------------------------------------------------------

\begin{document}

\selectthesislanguage

% Külső borító, minta kötéshez - csak elektronikus leadás esetén eltávolítandó
%~~~~~~~~~~~~~~~~~~~~~~~~~~~~~~~~~~~~~~~~~~~~~~~~~~~~~~~~~~~~~~~~~~~~~~~~~~~~~~~~~~~~~~
%\hypersetup{pageanchor=false}
%--------------------------------------------------------------------------------------
%	The outer cover page
%--------------------------------------------------------------------------------------
\pagenumbering{gobble}
\begin{center}

\vspace{130pt} %because it's the top
{\Large \bfseries \sze \\ \givk \\ \tanszek }\\
\vspace{180pt}
{\huge \bfseries \MakeUppercase {\doktipus}}\\
\vspace{160pt}
{\huge \bfseries{\szerzo}}

\vspace{40pt}
\Large \textbf{\szak{}}\\
\textbf{\szakirany}\\
\vfill
{\Large \textbf{\the\year}}

\end{center}
\hypersetup{pageanchor=false}



% Címoldal 
%~~~~~~~~~~~~~~~~~~~~~~~~~~~~~~~~~~~~~~~~~~~~~~~~~~~~~~~~~~~~~~~~~~~~~~~~~~~~~~~~~~~~~~
\hypersetup{pageanchor=false}
%--------------------------------------------------------------------------------------
%	The title page
%--------------------------------------------------------------------------------------
\begin{titlepage}

\ifthenelse{\equal{\tanszek}{\szeit}}{
    \includegraphics[height=15mm,keepaspectratio]{figures/sze_logo_wide.pdf} 
    \hfill
    \includegraphics[height=15mm,keepaspectratio]{figures/dept_it_logo.png}
}{
    \includegraphics[width=110mm,keepaspectratio]{figures/sze_logo_wide.pdf}
}

\begin{center}

\vspace{130pt} %because it's the top
{\Huge \bfseries \MakeUppercase {\szeGyor}}\\
\vspace{80pt}
{\huge \bfseries \cim}\\
%\vspace{10pt}
%{\Huge \bfseries {\doktipus}}\\
\vspace{80pt}
{\huge \bfseries{\szerzo}}
 
\vspace{80pt}
\Large \textbf{\szak{}}\\
\textbf{\szakirany}\\
\vfill
{\Large \textbf{\the\year}}

\end{center}
\end{titlepage}
\hypersetup{pageanchor=false}



%TODO Feladatkiíró lap helye, csak a nyomtatott verzióba kerül az eredeti példány
%~~~~~~~~~~~~~~~~~~~~~~~~~~~~~~~~~~~~~~~~~~~~~~~~~~~~~~~~~~~~~~~~~~~~~~~~~~~~~~~~~~~~~~
\pagenumbering{gobble}
%--------------------------------------------------------------------------------------
% Feladatkiiras (a tanszeken atveheto, kinyomtatott valtozat)
%--------------------------------------------------------------------------------------
\clearpage

% PDF formátumú leírás esetén
%\includepdf{figures/start.pdf}

% Képfájlokhoz
%\includegraphics*[width=\linewidth]{figures/start.png}

% Nyilatkozat és Kivonat
%~~~~~~~~~~~~~~~~~~~~~~~~~~~~~~~~~~~~~~~~~~~~~~~~~~~~~~~~~~~~~~~~~~~~~~~~~~~~~~~~~~~~~~
\masfelessorkoz{
% Tartalomjegyzék
%~~~~~~~~~~~~~~~~~~~~~~~~~~~~~~~~~~~~~~~~~~~~~~~~~~~~~~~~~~~~~~~~~~~~~~~~~~~~~~~~~~~~~~
\tableofcontents\vfill

% A dolgozat lényegi része
%~~~~~~~~~~~~~~~~~~~~~~~~~~~~~~~~~~~~~~~~~~~~~~~~~~~~~~~~~~~~~~~~~~~~~~~~~~~~~~~~~~~~~~
\newpage
\pagenumbering{arabic}
\setcounter{page}{1}

% Saját munka

%----------------------------------------------------------------------------
\chapter{\bevezetes}
%----------------------------------------------------------------------------

%----------------------------------------------------------------------------
\section{Bevezetés}
%----------------------------------------------------------------------------

%----------------------------------------------------------------------------
\chapter{\AudioBasics}
%----------------------------------------------------------------------------
%----------------------------------------------------------------------------
\chapter{\AudioOverIp}
%----------------------------------------------------------------------------
\section{Bevezetés az Audio over IP világába~\cite{AHNERT2023}}
%----------------------------------------------------------------------------
%----------------------------------------------------------------------------
\chapter{\DevelopementPlan}
%----------------------------------------------------------------------------
\section{Fejlesztés célja és kiindulópontja}
%----------------------------------------------------------------------------

A kompetenciaelemzés során világossá vált számomra, hogy vezetőként már most is számos erős területtel rendelkezem, ugyanakkor a fejlődés lehetősége folyamatos.
Célom, hogy tudatosan továbbfejlesszem azokat a készségeimet, amelyek a mindennapi vezetői munkám során még hatékonyabbá, kiegyensúlyozottabbá és inspirálóbbá tehetnek.
A fejlesztési terv alapját az önértékelő teszt eredményei adják, valamint a saját tapasztalataim és visszajelzéseim a kollégáimtól.

Kiemelten szeretnék foglalkozni három fő fejlesztési területtel:

\begin{enumerate}
    \item A stresszkezelés és energiamenedzsment tudatosabb egyensúlyba hozása.
    \item A jutalmazás és elismerés strukturáltabb, tudatosabb alkalmazása a csapatomban.
    \item Az önfejlesztés és inspiráció hosszú távú fenntartása, hogy vezetőként ne csak irányítsak, hanem példát is mutassak.
\end{enumerate}

%----------------------------------------------------------------------------
\section{Fejlesztési területek és célkitűzések}
%----------------------------------------------------------------------------

\subsection{Stresszkezelés és energiamenedzsment}

\textbf{Kiinduló helyzet:} \\
A munkám során gyakran kezelek több, párhuzamosan futó projektet, amelyekhez magas szintű koncentráció és folyamatos felelősségvállalás társul.
Noha általában jól reagálok stresszes helyzetekre, időnként hajlamos vagyok túlvállalni magam, ami hosszú távon mentális fáradtsághoz vezethet.

\textbf{Fejlesztési cél:} \\
Tudatos időbeosztás és energiamenedzsment kialakítása annak érdekében, hogy a hatékonyság mellett a regenerálódásra is elegendő idő jusson.

\textbf{Fejlesztési lépések:}
\begin{itemize}
    \item A heti munkabeosztás átnézése és a „pufferidők” beépítése a határidők közé.
    \item Rendszeres pihenőidők beiktatása (pl. heti 1 nap teljes offline nap).
    \item Relaxációs és koncentrációs technikák (pl. légzőgyakorlatok, rövid meditáció) tudatos használata.
    \item A sport és fizikai aktivitás beépítése a heti rutinba – nemcsak fizikai, hanem mentális feltöltődés céljából.
\end{itemize}

\textbf{Eredményességi mutató:} \\
Érezhető energiaszint-növekedés, kevesebb túlhajszoltság, kiegyensúlyozottabb döntéshozatal a projektek során.

\subsection{Jutalmazás és elismerés a csapatban}

\textbf{Kiinduló helyzet:} \\
A vezetői munkám során elsősorban személyes példamutatással és szakmai támogatással motiválom a kollégáimat.
Bár ez hatékony, úgy érzem, hogy a strukturáltabb visszajelzés- és elismerésrendszer erősíthetné a csapatkohéziót és az egyéni teljesítményt.

\textbf{Fejlesztési cél:} \\
Olyan, tudatosan felépített elismerési rendszer kialakítása, amely nemcsak szakmai, hanem emberi szinten is megerősíti a csapattagokat.

\textbf{Fejlesztési lépések:}
\begin{itemize}
    \item Havi rendszerességű személyes visszajelző beszélgetések bevezetése.
    \item Egyéni és csoportszintű teljesítmények elismerése szóban és írásban.
    \item Kis léptékű, de rendszeres jutalmazási formák (pl. közös program, kiemelés meetingeken).
    \item A csapattagok bevonása a motivációs eszközök közös megtervezésébe.
\end{itemize}

\textbf{Eredményességi mutató:} \\
Növekvő csapatmotiváció, pozitívabb visszajelzések a munkahelyi légkörről, alacsonyabb fluktuáció és jobb együttműködés.

\subsection{Önismeret és inspiráció fenntartása}

\textbf{Kiinduló helyzet:} \\
Vezetőként fontos számomra, hogy folyamatosan fejlődjek, nyitott maradjak az új megoldásokra, és ne ragadjak bele a rutinba.
Tudom, hogy a hosszú távú inspiráció és önfejlesztés kulcsa az, ha én magam is folyamatosan tanulok, fejlődök és új impulzusokat keresek.

\textbf{Fejlesztési cél:} \\
Az önismeret és a szakmai fejlődés tudatos fenntartása, valamint az inspiráció átadása a csapatom számára.

\textbf{Fejlesztési lépések:}
\begin{itemize}
    \item Rendszeres önreflexiós naplóvezetés (havonta egyszer).
    \item Évente legalább két szakmai konferencia vagy tréning látogatása.
    \item Inspiráló szakirodalmak, vezetéselméleti könyvek rendszeres olvasása.
    \item Mentori kapcsolat kialakítása vagy szakmai tapasztalatcsere egy hasonló területen dolgozó vezetővel.
\end{itemize}

\textbf{Eredményességi mutató:} \\
Tartós szakmai fejlődés, fokozódó motiváció, valamint a csapat inspirálása és fejlődése a saját példámon keresztül.

%----------------------------------------------------------------------------
\section{Összegzés}
%----------------------------------------------------------------------------

A fejlesztési tervem célja, hogy a vezetői kompetenciáimat a következő években tudatosan továbbfejlesszem, és stabil alapokra helyezzem a hosszú távú szakmai hatékonyságomat.
Úgy érzem, a legfontosabb lépés az, hogy megtartsam azt az önismeretet és felelősségtudatot, amely eddig is jellemezte a munkámat, miközben új szintre emelem a csapatom motiválását és a saját energiamenedzsmentemet.

Meggyőződésem, hogy a vezetés nem pozíció, hanem szemlélet — és a fejlődéshez való nyitottság az, ami igazán jó vezetővé teheti az embert.


% Ábrák listája - a word-ös sablon szerint szükséges
%~~~~~~~~~~~~~~~~~~~~~~~~~~~~~~~~~~~~~~~~~~~~~~~~~~~~~~~~~~~~~~~~~~~~~~~~~~~~~~~~~~~~~~
\listoffigures\addcontentsline{toc}{chapter}{\listfigurename}

% Táblázatok listája - opcionális
%~~~~~~~~~~~~~~~~~~~~~~~~~~~~~~~~~~~~~~~~~~~~~~~~~~~~~~~~~~~~~~~~~~~~~~~~~~~~~~~~~~~~~~
\listoftables\addcontentsline{toc}{chapter}{\listtablename}

% Irodalomjegyzék
%~~~~~~~~~~~~~~~~~~~~~~~~~~~~~~~~~~~~~~~~~~~~~~~~~~~~~~~~~~~~~~~~~~~~~~~~~~~~~~~~~~~~~~
\addcontentsline{toc}{chapter}{\bibname}
\bibliography{bib/references}

% Függelékek
%~~~~~~~~~~~~~~~~~~~~~~~~~~~~~~~~~~~~~~~~~~~~~~~~~~~~~~~~~~~~~~~~~~~~~~~~~~~~~~~~~~~~~~
%%----------------------------------------------------------------------------
\appendix
%----------------------------------------------------------------------------
\chapter*{\fuggelek}\addcontentsline{toc}{chapter}{\fuggelek}
\setcounter{chapter}{\appendixnumber}
%\setcounter{equation}{0} % a fofejezet-szamlalo az angol ABC 6. betuje (F) lesz
\numberwithin{equation}{section}
\numberwithin{figure}{section}
\numberwithin{lstlisting}{section}
%\numberwithin{tabular}{section}

%\label{page:last}
}
\end{document}